\section{Preliminary}\label{h1:prelim}

\paragraph{Notations.}
We use the notation $\simplex^d$ to denote the $(d-1)$-dimensional standard simplex in $\dR^d$: $\simplex^d = \cbrm[\big]{x\in \dR^d_+ \mid  \sum_{i=1}^d x_i=1}$. We similarly use $\simplex^A$ for a finite set $A$ to denote the set of distributions over $A$. %
A sequence of vectors changing over time is indexed using superscript for the time dimension and subscript for the inherent dimension of the vector; for example, $v!t_i$ is the $i$-th dimension of the vector $v$ at the $t$-th round. %
For an integer $m$, $[m]$ represents the set $\{1, \dots, m\}$.

\paragraph{Problem Statement.}
A learner plays a zero-sum game against an adversary. They can choose their actions from finite action sets $\calX\subset \dR^{d_x}$ and $\calY\subset \dR^{d_y}$ respectively, which have cardinalities $m_x=|\calX|$ and $m_y=|\calY|$.
The game is identified with a utility function $u : \calX\times\calY\to  [-1, 1]$, where $u(x,y)$ is the expected reward received by the learner if she plays $x$ and the adversary plays $y$.
We assume the utility function is bilinear, i.e., $u(x,y)=\inner{x}{Ay}$ for some unknown matrix $A\in \dR^{d_x\times d_y}$.
This general formulation captures all finite normal-form games as special cases where $\calX$ and $\calY$ are standard basis sets.

The learner obtains information about the game by repeatedly playing.
Specifically, in each round $t=1, 2, \dots$, the learner selects an action $x!t \in \calX$, and the adversary simultaneously selects $y!t \in \calY$. We consider an \emph{adaptive adversary} who can select $y!t$ based on the realization of the history, including the learner's past actions
$x!1, \dots, x!{t-1}$.
After both players commit to their actions, the player receives a noisy reward $r_t = u(x!t, y!t) + \eps_t$, assumed to be in $[-1, 1]$. 
The noise sequence $\{\eps_t\}$ is a martingale difference sequence; specifically, conditional on the history and chosen actions, the bias $\EE[\eps_t]$ is $0$. %

\paragraph{Feedback Models.}
We study two feedback settings depending on the information the learner observes after each round:
\begin{itemize}[noitemsep,topsep=3pt]
    \item {\itshape Uninformed}: The learner only observes  the reward $r_t$. This is the same feedback model used in the self-play analysis by \citet{ito2025instancedependent}.
    \item {\itshape Informed}: The learner observes both the reward $r_t$ and the adversary's action $y!t$. This is the feedback model used by \citet{odonoghue2021matrix} and \citet{maiti2025limitations}.
\end{itemize}


\paragraph{Properties of the Game.}
If a pair of actions $(x^*,y^*)\in \calX\times\calY$ satisfies \begin{equation}
    u(x,y^*)\leq u(x^*,y^*)\leq u(x^*,y), \label{eq:ne}
\end{equation}
for any $x\in \calX$ and $y\in \calY$, we say that $(x^*,y^*)$ is a \emph{pure-strategy Nash equilibrium} (PSNE). If \eqref{eq:ne} holds with strict inequality whenever $x\neq x^*$ and $y\neq y^*$, we say the pair is a \emph{strict} PSNE.

For a pair of distributions over actions $(p, q)\in \simplex^\calX\times \simplex^\calY$, we slightly abuse the notation of the utility function and define $u(p,q)=\EE_{x\sim p,y\sim q}[u(x,y)]=\sum_{x\in \calX}\sum_{y\in \calY}p_x q_y u(x,y)$ to be the expected reward of the learner if she plays a sample from $p$ and the adversary plays a sample from $q$. If a pair of distributions over actions $(p^*, q^*)\in \simplex^\calX\times \simplex^\calY$ satisfies
\begin{equation}
    u(p,q^*)\leq u(p^*,q^*)\leq u(p^*,q),
\end{equation}
for any $p\in \simplex^\calX$ and $q\in \simplex^\calY$, we say that $(p^*,q^*)$ is a \emph{mixed-strategy Nash equilibrium} (MSNE). The celebrated minimax theorem by \citet{vonneumann1928theory} proves that an MSNE always exists, and that all Nash equilibria yield the same expected value $\vMix=u(p^*,q^*)$, which is also equal to $\max_{p\in \simplex^\calX}\min_{q\in \simplex^\calY} u(p, q)$ and $\min_{q\in \simplex^\calY}\max_{p\in \simplex^\calX} u(p, q)$.

In contrast, central to our work is the concept of pure-strategy maximin value, defined as $\vStar\defeq \max_{x\in \calX} \min_{y\in \calY} u(x,y)$, which is the maximum utility a deterministic learner can guarantee in the worst case.
By definition, $\vStar$ is no more than $\vMix$, but when a PSNE exists, these two values coincide.

\paragraph{Goal.}
From the learner's perspective, this problem can be seen as an instance of adversarial multi-armed bandits~\citep{auer2002nonstochastic}, where the standard goal is to minimize \textit{external regret}, which compares the learner's total reward to that of the best fixed action in hindsight.
Specifically, let $\ER_T(x) \defeq \EE\sbrm[\Big]{\sum_{t=1}^T \rbrm[\big]{u(x,y!t)-u(x!t,y!t)}}$ be the regret compared to action $x$.
Then external regret is defined as
$\ER_T \defeq \max_{x\in \calX} \cbrm[\big]{\ER_T(x)}$.
Applying standard algorithms such as Exp3~\citep{auer2002nonstochastic} guarantees $\ER_T = \tilde{\order}(\sqrt{m_xT})$, which is generally not improvable even for simple games, as shown in~\citet[Theorem~3]{maiti2025limitations}.

To get around this obstacle,
we propose to consider the following new regret metric, termed \emph{Pure-Strategy Maximin Regret} (PSMR), as the learner's performance measure:
\begin{equation}
    \psmr_T\defeq \EE\sbrm[\Bigg]{\sum_{t=1}^T \rbrm[\big]{\vStar-u(x!t,y!t)}},
    \label{eq:psmr-def}
\end{equation}
where the expectation is taken over the randomness of the algorithm, the noise in the reward, and potentially the adversary.
In words, PSMR measures the learner's accumulated deficit compared to the pure-strategy maximin value.
On the other hand, the Nash-value regret considered in prior work such as~\citet{odonoghue2021matrix, maiti2025limitations} is defined as $\NR_T \defeq\EE\sbrm[\Big]{\sum_{t=1}^T \rbrm[\big]{\vMix-u(x!t,y!t)}}$, which measures the learner's accumulated deficit compared to the Nash value (equivalently, the mixed-strategy maximin value).

By definition, it is clear that $\psmr_T \leq \NR_T \leq \ER_T$; the first inequality is due to $\vStar\leq \vMix$, and the second due to $\NR_T=\EE\sbrm[\Big]{\sum_{t=1}^T \rbrm[\big]{u(p^*, q^*)-u(x!t,y!t)}}\leq\EE_{x\sim p^*}[\ER_T(x)]\leq \ER_T$ for any NE $(p^*,q^*)$.
When the game has a PSNE, however, $\psmr_T$ and $\NR_T$ coincide. 

Finally, we point out that in the special case when the opponent has only one action ($|\calY|=1$), our problem is simply the standard stochastic linear bandit~\citep{abbasi-yadkori2011improved} or
stochastic multi-armed bandit~\citep{lai1985asymptotically} if $\calX$ is the standard basis.
In this case,
both $\psmr_T$ and $\NR_T$ recover $\ER_T$.
Therefore, our results can also be seen as a generalization of the classical instance-dependent bounds for stochastic bandits achieved by, for example, the UCB algorithm~\citep{auer2002using}.


